\chapter{Thiết kế Hệ thống}
\label{chap:design}

\section{Kiến trúc Tổng thể}
\label{sec:design-architecture}

Ứng dụng được tổ chức theo kiến trúc \textbf{3 lớp} (Three-Tier Architecture):

\begin{enumerate}
  \item \textbf{Lớp Trình bày (Presentation Layer)}: Các màn hình React Native
        trong \texttt{src/app/} và component UI trong \texttt{src/components/}.
        Lớp này chỉ có nhiệm vụ hiển thị kết quả và thu thập đầu vào từ người dùng.
  \item \textbf{Lớp Nghiệp vụ (Business Logic Layer)}: Các hàm tính toán thuần túy
        (pure functions) trong \texttt{src/lib/} và \texttt{src/services/motor-dk-selection.ts}.
        Lớp này hoàn toàn độc lập với UI và cơ sở dữ liệu,
        cho phép kiểm thử đơn vị dễ dàng.
  \item \textbf{Lớp Dữ liệu (Data Layer)}: Drizzle ORM + Expo SQLite
        (\texttt{src/db/}) cho dữ liệu cục bộ;
        Axios client (\texttt{src/services/api/}) cho giao tiếp với backend Java Spring.
\end{enumerate}

Luồng dữ liệu chính giữa 3 module tính toán:
màn hình \texttt{calc.tsx} (Module~1) gọi các hàm nghiệp vụ,
nhận kết quả rồi ghi vào \texttt{globalNavigationState};
khi điều hướng đến \texttt{chain.tsx} (Module~2) hoặc \texttt{gear.tsx} (Module~3),
màn hình đọc state và tự động hiển thị kết quả mà không cần người dùng nhập lại.

\section{Áp dụng Design Pattern (Module 5 -- CNPM)}
\label{sec:design-patterns}

\subsection{Strategy Pattern -- Xử lý Công thức Tính toán}

\textbf{Bài toán}: Bộ truyền bánh răng có hai biến thể khác nhau về công thức
(bánh răng côn cấp nhanh và bánh răng trụ cấp chậm).
Nếu viết cứng cả hai vào một hàm duy nhất, code sẽ phức tạp và khó bảo trì.
Ngoài ra, bộ truyền xích cần cho phép thay thế các hằng số tính toán
(hệ số môi trường, hệ số sử dụng) mà không sửa logic chính.

\textbf{Giải pháp}: Mỗi biến thể được triển khai thành một hàm tính toán độc lập
(Concrete Strategy) với kiểu đầu vào và đầu ra được định nghĩa rõ ràng:

\begin{verbatim}
// Concrete Strategy 1: Bánh răng côn (cấp nhanh)
calculateBevelGearStage(input: BevelGearInput): BevelGearResult

// Concrete Strategy 2: Bánh răng trụ (cấp chậm)
calculateCylindricalGearStage(input: CylindricalGearInput): CylindricalGearResult

// Strategy với tham số cấu hình (bộ truyền xích)
calculateChainDrive(input: ChainDriveInput,
                   constants: ChainDriveConstants = DEFAULT_CONSTANTS): ChainDriveResult
\end{verbatim}

Màn hình \texttt{gear.tsx} đóng vai trò Context: nó gọi lần lượt
\texttt{calculateBevelGearStage} và \texttt{calculateCylindricalGearStage}
theo cùng một giao diện thống nhất, dễ dàng mở rộng thêm biến thể mới mà không
ảnh hưởng đến phần còn lại.

\subsection{Observer Pattern -- Đồng bộ Dữ liệu Giữa các Module}

\textbf{Bài toán}: Khi người dùng hoàn thành Module~1,
Module~2 và Module~3 phải tự động nhận các tham số tương ứng
($P_{III}$, $n_{III}$, $T_I$, $T_{II}$, \ldots) mà không yêu cầu người dùng nhập lại.

\textbf{Giải pháp}: Đối tượng \texttt{globalNavigationState}
(định nghĩa trong \texttt{src/lib/globalState.ts}) đóng vai trò Subject (Observable).
Khi màn hình \texttt{calc.tsx} hoàn thành Module~1, nó cập nhật:

\begin{verbatim}
// Subject cập nhật trạng thái
globalNavigationState.gearResult = {
  bevelResult: calculateBevelGearStage(...),
  cylindricalResult: calculateCylindricalGearStage(...)
};
globalNavigationState.chainResult = calculateChainDrive(...);
\end{verbatim}

Các màn hình \texttt{gear.tsx} và \texttt{chain.tsx} (đóng vai trò Observer)
đọc state này trong hook \texttt{useFocusEffect} khi màn hình được kích hoạt,
tự động nạp và hiển thị kết quả đã tính.

\subsection{Facade Pattern -- Giao diện Thống nhất cho Logic Tính toán}

\textbf{Bài toán}: Mỗi module tính toán gồm nhiều bước phức tạp:
tra bảng, tính trung gian, kiểm nghiệm nhiều điều kiện.
Màn hình (UI) không nên biết chi tiết từng bước nội bộ
-- điều đó làm tăng coupling và khó kiểm thử.

\textbf{Giải pháp}: Mỗi module cung cấp một hàm Facade duy nhất.
Toàn bộ quy trình nhiều bước được gói gọn bên trong:

\begin{verbatim}
// Facade Module 2: 1 lời gọi → toàn bộ kết quả xích
const result = calculateChainDrive({ P, n, u });
// result.powerParams  — thông số công suất và hệ số
// result.geometry     — z1, z2, pc, a, X, L, v, Ft, ...
// result.validations  — 4 kết quả kiểm nghiệm (pass/fail)

// Facade Module 3 cấp nhanh: 1 lời gọi → 7 bước tính
const bevel = calculateBevelGearStage({ T_I, n_I, u_1, L_h });
// bevel.step1 .. bevel.step7  — từng bước trung gian
// bevel.sigma_H, sigma_F1, sigma_F2  — kết quả kiểm nghiệm
\end{verbatim}

Màn hình chỉ cần gọi một hàm và nhận cấu trúc kết quả đầy đủ để hiển thị,
không cần biết có bao nhiêu bước trung gian hay bảng tra cứu nào được dùng.

\subsection{Singleton Pattern -- Quản lý Cơ sở Dữ liệu Cục bộ}

\textbf{Bài toán}: Nhiều module (seed, truy vấn động cơ, lưu phiên)
đều cần truy cập cơ sở dữ liệu SQLite.
Nếu mỗi module mở một kết nối riêng sẽ gây xung đột file lock và lãng phí bộ nhớ.

\textbf{Giải pháp}: File \texttt{src/db/client.ts} xuất một instance Drizzle duy nhất:

\begin{verbatim}
// src/db/client.ts
import { drizzle } from 'drizzle-orm/expo-sqlite';
import { openDatabaseSync } from 'expo-sqlite';

const sqlite = openDatabaseSync('mechcalc.db');
export const db = drizzle(sqlite);   // Singleton toàn cục
\end{verbatim}

Trong JavaScript/TypeScript, module-level exports được khởi tạo đúng một lần
nhờ cơ chế module caching của Node.js/Metro bundler,
đảm bảo tính chất Singleton mà không cần class hay static method.
Mọi file \texttt{import \{ db \} from '@/db/client'} đều nhận cùng một instance.

\section{Thiết kế Cơ sở Dữ liệu}
\label{sec:design-db}

Cơ sở dữ liệu cục bộ gồm 2 bảng catalog động cơ tiêu chuẩn.
Dữ liệu phiên tính toán được lưu trên backend PostgreSQL.

\subsection{Bảng \texttt{motors\_dk} -- Catalog Động cơ Series ĐK/IE3}

\begin{table}[H]
  \centering
  \caption{Schema bảng \texttt{motors\_dk}}
  \label{tab:schema-motors-dk}
  \begin{tabular}{llll}
    \toprule
    \textbf{Cột} & \textbf{Kiểu} & \textbf{Ràng buộc} & \textbf{Mô tả} \\
    \midrule
    id                    & INTEGER & PK, AUTOINCREMENT & Khóa chính \\
    model                 & TEXT    & NOT NULL & Mã hiệu (vd: IE3-W41R 160 M2) \\
    power\_kw             & REAL    & NOT NULL & Công suất định mức (kW) \\
    speed\_rpm            & INTEGER & NOT NULL & Tốc độ quay định mức (v/p) \\
    cos\_phi              & REAL    & NOT NULL & Hệ số công suất \\
    starting\_torque\_ratio & REAL  & NOT NULL & Tỉ số mô-men mở máy $M_a/M_b$ \\
    max\_torque\_ratio    & REAL    & NOT NULL & Tỉ số mô-men cực đại $M_k/M_b$ \\
    rotor\_inertia\_gd2   & REAL    & NOT NULL & Mô-men quán tính rotor (kg$\cdot$m$^2$) \\
    weight\_kg            & REAL    & NOT NULL & Khối lượng (kg) \\
    poles                 & INTEGER & NOT NULL & Số cực \\
    sync\_speed\_rpm      & INTEGER & NOT NULL & Tốc độ đồng bộ (v/p) \\
    \bottomrule
  \end{tabular}
\end{table}

Bảng \texttt{motors\_4a} có cấu trúc tương tự, lưu catalog động cơ series 4A.
Cả hai được nạp tự động từ seed data TypeScript khi ứng dụng khởi động lần đầu.

\section{Thiết kế Giao diện Người dùng (UI/UX)}
\label{sec:design-ui}

\subsection{Nguyên tắc Thiết kế}

Giao diện được thiết kế theo các nguyên tắc sau:
\begin{itemize}
  \item \textbf{Mobile-first}: Tối ưu cho màn hình điện thoại dọc (portrait mode).
  \item \textbf{Minimal input}: Tự động chuyển tham số liên module,
        giảm tối đa số lần người dùng phải nhập thủ công.
  \item \textbf{Step-by-step display}: Kết quả tính toán được trình bày
        theo từng bước có đánh số,
        giúp sinh viên đối chiếu trực tiếp với tính tay theo giáo trình.
  \item \textbf{Visual feedback}: Kết quả kiểm nghiệm đạt/không đạt
        hiển thị bằng màu xanh lá/đỏ kèm giá trị số rõ ràng.
\end{itemize}

\subsection{Luồng Màn hình (Screen Flow)}

Luồng điều hướng chính của ứng dụng được mô tả như sau:

\begin{enumerate}
  \item \textbf{Xác thực} (\texttt{(auth)/login.tsx}, \texttt{register.tsx}):
        Người dùng đăng nhập hoặc đăng ký tài khoản.
        \texttt{AuthGuard} trong root layout tự động chuyển hướng nếu chưa đăng nhập.
  \item \textbf{Trang chủ} (\texttt{(tabs)/index.tsx}):
        Chào mừng, truy cập nhanh vào Tính toán và Dự án.
  \item \textbf{Module~1 -- Tính động cơ} (\texttt{(tabs)/calc.tsx}):
        Nhập $P_{lv}$, $n_{lv}$, $L$ $\to$ hiển thị từng bước tính
        $\to$ chọn động cơ từ danh sách gợi ý
        $\to$ xem bảng phân phối $P$--$n$--$T$ trên 4 trục.
  \item \textbf{Module~2 -- Tính xích} (\texttt{chain.tsx}):
        Tự động nạp tham số từ Module~1
        $\to$ hiển thị thông số bộ truyền xích và 4 điều kiện kiểm nghiệm.
  \item \textbf{Module~3 -- Tính bánh răng} (\texttt{gear.tsx}):
        Tự động nạp tham số từ Module~1
        $\to$ hiển thị kết quả cấp nhanh (bánh răng côn)
        và cấp chậm (bánh răng trụ).
  \item \textbf{Quản lý Dự án} (\texttt{(tabs)/project.tsx}, \texttt{project/[id].tsx}):
        Danh sách phiên tính toán đã lưu; xem chi tiết từng phiên.
  \item \textbf{Trang cá nhân} (\texttt{(tabs)/profile.tsx}):
        Thông tin tài khoản, đăng xuất.
\end{enumerate}
